% ---
% id: receiver/lem:fock-maps
% kind: lem
% title: "The lift of maps between mode spaces and of direct sums"
% summary: >-
%   The Fock lift respects products and adjoints of maps between two mode
%   spaces and takes contractions to contractions and isometries to
%   isometries. On disjoint mode sets the occupation identification is
%   unitary, and the lift of a direct sum is the tensor product of the lifts.
% prominence: technical
% status: accepted
% uses: [def:fock, lem:fock-functor]
% assumes: []
% refs: []
% provenance: newly-derived
% lean:
% ---

Lemma~\ref{lem:fock-functor} is stated for operators on one space. A loss in
front of a receiver is an isometry into a larger mode space, so the same
properties are needed for maps between two spaces, together with the
identification of the Fock space over two disjoint mode sets with the tensor
product of their Fock spaces.

\begin{lemma}\label{lem:fock-maps}
Let $\mathcal{K}_0,\mathcal{K}_1,\mathcal{K}_2$ be finite-dimensional, and let
$X\colon\mathcal{K}_0\to\mathcal{K}_1$ and $Y\colon\mathcal{K}_1\to\mathcal{K}_2$
be linear maps.
\begin{enumerate}[label=(\roman*),leftmargin=3.0em]
\item $\Fock(YX)=\Fock(Y)\Fock(X)$ and $\Fock(X)^{\dagger}=\Fock(X^{\dagger})$;
  if $X^{\dagger}X\preceq\id$ then $\lVert\Fock(X)\rVert\le1$, and if
  $X^{\dagger}X=\id$ then $\Fock(X)$ is an isometry.
\item Let $S$ and $T$ be disjoint finite mode sets. The linear map
  $\mathcal{F}(\mathbb{C}^{S})\otimes\mathcal{F}(\mathbb{C}^{T})\to
  \mathcal{F}(\mathbb{C}^{S\sqcup T})$ defined on occupation vectors
  (Definition~\ref{def:fock}) by
  $\ket{\vec m_1}\otimes\ket{\vec m_2}\mapsto\ket{\vec m_1,\vec m_2}$, the
  \emph{occupation identification},
  is unitary, and on its $n$-photon level it is
  $\ket{\varphi}\otimes\ket{\chi}\mapsto\binom{n}{p}^{1/2}P_n(\varphi\otimes\chi)$
  for $\varphi\in\Symn{p}(\mathbb{C}^{S})$, $\chi\in\Symn{q}(\mathbb{C}^{T})$,
  $p+q=n$, with $P_n$ the orthogonal projector of
  $(\mathbb{C}^{S\sqcup T})^{\otimes n}$ onto its symmetric subspace.
\item Under that identification, for operators $X_1$ on $\mathbb{C}^{S}$ and
  $X_2$ on $\mathbb{C}^{T}$,
  $\Fock(X_1\oplus X_2)=\Fock(X_1)\otimes\Fock(X_2)$.
\end{enumerate}
\end{lemma}

\begin{proof}
(i) $X^{\otimes n}$ maps $\Symn{n}(\mathcal{K}_0)$ into
$\Symn{n}(\mathcal{K}_1)$ (Definition~\ref{def:fock}), and so does
$(X^{\dagger})^{\otimes n}$ from $\Symn{n}(\mathcal{K}_1)$ back, both
intertwining the permutation actions. Then $(YX)^{\otimes n}=Y^{\otimes n}X^{\otimes n}$
restricted to $\Symn{n}(\mathcal{K}_0)$ is the product of the restrictions,
and the restriction of $(X^{\otimes n})^{\dagger}=(X^{\dagger})^{\otimes n}$ to
$\Symn{n}(\mathcal{K}_1)$ is the adjoint of the restriction of
$X^{\otimes n}$, since both restrictions map between the invariant subspaces;
summing over $n$ gives the two identities, as in
Lemma~\ref{lem:fock-functor}\,(ii). If $X^{\dagger}X\preceq\id$ then
$\lVert X\rVert\le1$, so every block $X^{\otimes n}$ has norm at most one, as
in Lemma~\ref{lem:fock-functor}\,(iii). If $X^{\dagger}X=\id$, the two
identities give $\Fock(X)^{\dagger}\Fock(X)=\Fock(X^{\dagger}X)=\Fock(\id)=\id$.

(ii) For an occupation vector $\vec m$ with $|\vec m|=n$ write
$e^{\otimes\vec m}$ for the tensor product of $m_1$ copies of the first basis
vector, then $m_2$ of the second, and so on, and $\vec m!\coloneqq\prod_jm_j!$.
$P_ne^{\otimes\vec m}$ is $\vec m!/n!$ times the sum of the $n!/\vec m!$
distinct reorderings of $e^{\otimes\vec m}$, which are orthonormal, so
$\lVert P_ne^{\otimes\vec m}\rVert^{2}=\vec m!/n!$ and the normalized
symmetrization of Definition~\ref{def:fock} is
$\ket{\vec m}=(n!/\vec m!)^{1/2}P_ne^{\otimes\vec m}$. Let $J_{p,q}$ be the
displayed map on the level $p+q=n$. Since $P_n(P_p\otimes P_q)=P_n$, the
product $P_p\otimes P_q$ averaging over a subgroup of the permutations, and a
reordering of the factors is a permutation,
\begin{equation*}
  J_{p,q}\big(\ket{\vec m_1}\otimes\ket{\vec m_2}\big)
  =\Big(\tbinom{n}{p}\tfrac{p!\,q!}{\vec m_1!\,\vec m_2!}\Big)^{1/2}
   P_n\big(e^{\otimes\vec m_1}\otimes e^{\otimes\vec m_2}\big)
  =\ket{\vec m_1,\vec m_2}.
\end{equation*}
So $\bigoplus_{p+q=n}J_{p,q}$ maps an orthonormal basis of the $n$-photon
level of the tensor product onto one of $\Symn{n}(\mathbb{C}^{S\sqcup T})$,
and the identification, their sum over $n$, is unitary.

(iii) $(X_1\oplus X_2)^{\otimes n}$ commutes with $P_n$
(Lemma~\ref{lem:fock-functor}\,(i)) and acts on
$(\mathbb{C}^{S})^{\otimes p}\otimes(\mathbb{C}^{T})^{\otimes q}$ as
$X_1^{\otimes p}\otimes X_2^{\otimes q}$, so
$(X_1\oplus X_2)^{\otimes n}J_{p,q}(\varphi\otimes\chi)
=J_{p,q}\big(X_1^{\otimes p}\varphi\otimes X_2^{\otimes q}\chi\big)$ for
$\varphi\in\Symn{p}(\mathbb{C}^{S})$ and $\chi\in\Symn{q}(\mathbb{C}^{T})$,
which is the claim on every level.
\end{proof}
